% Revised Conclusion and Future Work with ontology insights (paste into your thesis)
\chapter{Conclusion and Future Work}

This thesis presents an end-to-end framework that integrates (i) robust representation learning for noisy, low-data histopathology and (ii) ontology-grounded explainability that links morphologic evidence to genetic hypotheses and treatment knowledge in lung adenocarcinoma.

\paragraph{Key Findings on Real Histopathology (Zenodo ANORAK).}
The primary empirical contribution is a real-data ablation demonstrating that \emph{loss-induced embedding geometry} is a first-order determinant of generalization under clinical noise, class imbalance, and limited sample size. In the benchmark on Zenodo ANORAK (six-class LUAD pattern classification), the best configuration---\textbf{FuzzyArcLoss V3 with sub-centers}---achieved \textbf{76.96\% macro-F1} and \textbf{76.07\% accuracy}. Across the same protocol, multiple margin-based alternatives (including ArcFace-family variants) underperformed this configuration, confirming that both (i) uncertainty-aware margin modulation and (ii) explicit modeling of intra-class multimodality are necessary to stabilize representation learning for heterogeneous LUAD morphologies.

A subsequent advanced benchmark explored strategies layered on top of the learned representations, including alternative backbones, multi-head architectures, and ensemble methods. None of these configurations surpassed the loss-ablation winner, indicating that performance is constrained by the epistemic structure of patch-level supervision itself---specifically, limited effective sample size for minority patterns, intrinsic morphologic overlap between classes, and the use of hard labels for biologically fuzzy boundaries.

\paragraph{From Histologic Pattern Learning to Genetic Mutation Inference.}
A second central finding is that, even when patch-level classification reaches a ceiling, \emph{aggregated morphologic profiles} preserve clinically meaningful molecular signal. By converting patch-level outputs into case-level features---including total tile count and predicted pattern proportions---the framework enables gene-specific mutation inference with transparent attribution via SHAP. The results reveal mutation-dependent morphologic drivers: TP53 predictions are primarily driven by micropapillary and solid-associated components, whereas EGFR predictions depend more strongly on tile sampling volume and papillary/lepidic/acinar composition. This contrast supports a clinically grounded interpretation: mutations associated with global architectural disruption are more amenable to inference from aggregated morphology than mutations linked to subtle epithelial differentiation.

\paragraph{Explainability Contribution.}
The framework introduces an ontology-grounded explanation layer that links morphologic and molecular outputs to standardized biomedical knowledge without influencing training or prediction. Rather than a single monolithic OWL parse, the system uses a \emph{three-tier architecture}: (i)~a \textbf{curated knowledge graph} built from canonical IRIs (NCIt for patterns, genes, drugs, and diagnoses; MONDO for disease cross-reference), with edges from WHO~2021, OncoKB/FDA, TCGA/CIViC, and COSMIC carrying explicit \texttt{provenance}; (ii)~a \textbf{dynamic case-specific subgraph} that retains only entities relevant to each patient---patterns above a composition threshold, genes reachable from those patterns or predicted by the mutation models, and their associated treatments---so that explanations reflect the actual case rather than a generic schema; and (iii)~a \textbf{literature-driven enrichment pipeline} (DeepSearch) that periodically queries PubMed, Semantic Scholar, and arXiv for LUAD-related literature, extracts relation triples via an LLM, performs entity linking to the same IRIs and SHACL-like validation, and merges validated relations into the graph as versioned snapshots with a changelog. SPARQL is used when a triplestore is populated; otherwise the same IRI-based curated graph is served directly. A natural-language explanation agent (LLM) summarizes the graph for the clinician, with guardrails to avoid definitive diagnostic language. This design keeps explanations auditable, reproducible, and faithful to the computed evidence: every edge is traceable to either curated sources (asserted) or literature and morphology--molecular correlations (inferred).

Overall, the inability to surpass the observed patch-level performance ceiling should not be interpreted as a limitation of optimization or model capacity, but as empirical confirmation that \emph{hard tile-level supervision is structurally misaligned with intrinsically continuous morphologic phenomena}. The most reliable gains arise from objectives that explicitly acknowledge uncertainty and heterogeneity, suggesting that future progress depends more on supervision design and contextual modeling than on incremental architectural complexity.

\section*{Future Work}
To exceed the empirically observed ceiling on real histopathology data, future work should focus on improved supervision and context integration:

\begin{itemize}
  \item \textbf{Soft or ordinal supervision for morphologic continua:} replace hard class labels with uncertainty-aware or ordinal representations that explicitly encode transitions between glandular patterns.
  \item \textbf{Agreement-aware training:} incorporate multi-expert annotation, consensus filtering, or disagreement-aware weighting to reduce irreducible label noise at the tile level.
  \item \textbf{Multi-instance and slide-level modeling (MIL/WSI):} incorporate whole-slide context to resolve local ambiguity using global architectural organization and spatial co-occurrence.
  \item \textbf{Multi-scale modeling:} jointly model multiple magnifications to capture both gland-level architecture and fine cytologic detail.
  \item \textbf{Calibration and imbalance-aware mutation learning:} improve rare-mutation inference (e.g., EGFR, STK11, NF1) via cost-sensitive objectives, threshold calibration, and prevalence-aware evaluation.
  \item \textbf{Prospective validation:} assess clinical utility, calibration stability, and explanation usability in real diagnostic workflows, including pathologist-facing interfaces.
  \item \textbf{Ontology and knowledge-graph extension:} integrate additional ontologies (e.g., DOID, SNOMED~CT, ICD-O) and align pattern/gene/drug IRIs across institutions; formalize relations in OWL and evaluate OWL reasoners for consistency and classification.
  \item \textbf{DeepSearch and literature enrichment:} scale the literature pipeline to more sources and queries; strengthen entity linking and SHACL validation; support human-in-the-loop review of discovered relations before they enter the case graph.
  \item \textbf{Explanation quality and evaluation:} define metrics for ontology-generated explanations (completeness, consistency with evidence, clinician usefulness) and compare LLM summaries against structured SPARQL-based explanations in user studies.
\end{itemize}
